\documentclass[11pt]{article}

\usepackage{fontspec}
\setmainfont{TeX Gyre Pagella}
\usepackage{amsmath}
\usepackage{amssymb}
\usepackage{amsthm}
\usepackage[margin=1in]{geometry}

\newtheorem{theorem}{Theorem}
\newtheorem{proposition}{Proposition}
\newtheorem{definition}{Definition}
\newtheorem{example}{Example}
\newtheorem{corollary}{Corollary}

\title{Verification Inheritance Under Transformation \\
\large Robust Commitments, Claim Correspondence, and Epistemic Validity in Mathematical Documents}
\author{Flyxion \\ \small Independent Researcher}
\date{August 2026}

\begin{document}

\maketitle

\begin{abstract}
When a scholarly artifact containing verified claims is transformed, reformatted, translated, paraphrased, or edited, existing verification results do not automatically transfer to the transformed artifact. This paper isolates the resulting correspondence problem from the surrounding infrastructure of watermarking, provenance metadata, and cryptographic attestation, and treats it as a decision problem in its own right. We show that correspondence, defined as mutual semantic entailment between a claim and its transformed successor, is undecidable in general for closed first-order sentences over any signature containing a relation symbol of arity at least two. We then exhibit a restricted but non-trivial class of transformations, generated by a terminating and confluent rewrite system, under which correspondence is decidable by normalization. The two results together yield a principled three-valued decision procedure that certifies preservation only within an explicitly bounded regime, and otherwise reports that a transformation falls outside that regime rather than asserting a semantic verdict it cannot support. This paper does not propose a new watermarking, provenance, or trust-management mechanism; it argues that any such mechanism must confront the correspondence problem explicitly, and it supplies the boundary that mechanism should respect.
\end{abstract}

\section{Introduction}

Verification is typically treated as a property of an artifact at the moment it is checked. A theorem is proved, a computation is checked, a claim is reviewed, and the result is recorded. But artifacts do not remain fixed. They are reformatted, translated, paraphrased, excerpted, and edited, and each transformation raises the same question: does an earlier verification result still apply to the transformed artifact?

The question is not new. Proof repair asks whether an existing formal proof survives a change to the underlying library or definitions on which it depends. Citation-treatment systems in law ask whether an authority still supports the proposition for which it was cited after subsequent developments. Both traditions establish that support does not persist automatically across change; it must be re-established or explicitly inherited. This paper treats the same structural problem as it arises for mixed formal and informal scholarly claims under arbitrary transformation, and asks specifically what a system can and cannot decide about it.

The temptation, when building infrastructure for this problem, is to reach directly for a watermark or a provenance record and treat persistence of that record as evidence of continued validity. This is a category error. A watermark, a cryptographic commitment, or a provenance chain can establish that a transformed artifact traces back to an earlier one. None of these mechanisms can, by itself, establish that the claims the earlier artifact made are still the claims the transformed artifact makes. The two questions are independent, and conflating them allows a robust but semantically stale pointer to authenticate a claim that no longer holds.

This paper isolates the second question and treats it formally. Section 2 fixes a claim language and defines correspondence between a claim and its transformed successor. Section 3 proves that correspondence, so defined, is undecidable in general. Section 4 defines a restricted class of transformations under which correspondence becomes decidable by normalization, and gives a non-trivial rewrite system exhibiting the fragment. Section 5 assembles the two results into a three-valued decision procedure and states the epistemic discipline it enforces. Section 6 situates the result relative to prior work in watermarking, provenance, trust management, proof repair, and legal citation treatment, none of which is superseded by this paper; each supplies infrastructure that this paper's boundary constrains rather than replaces.

This paper does not attempt to build or evaluate a system. A companion paper takes the boundary established here as a design constraint and specifies a protocol, comprising a robust in-band locator, a signed claim manifest, transformation attestations, and a claim-type-indexed trust policy, that respects it. A further companion paper subjects an actual mathematical document to a battery of transformations, chosen to include both the certified fragment of this paper and adversarial cases deliberately outside it, and measures locator survival, claim correspondence, and verification inheritance as three independent quantities. Neither companion is required for the results below to hold.

\section{Preliminaries}

Fix a finite relational signature $\Sigma = \{P, Q, S\}$, where $P$ and $Q$ are unary relation symbols and $S$ is a binary relation symbol. The presence of a relation symbol of arity at least two is essential and is discussed after Theorem 1. A \emph{claim} is a closed sentence of first-order logic over $\Sigma$, built from atomic formulas using $\neg, \wedge, \vee, \rightarrow, \forall, \exists$ in the usual way, with no free variables.

Semantic entailment $c \models c'$ has its standard meaning: every $\Sigma$-structure satisfying $c$ also satisfies $c'$. A sentence is \emph{valid} if it is satisfied by every $\Sigma$-structure.

\begin{definition}[Correspondence]
For claims $c, c'$, define
\[
C(c,c') \iff (c \models c') \wedge (c' \models c).
\]
\end{definition}

$C(c,c')$ holds exactly when $c$ and $c'$ are logically equivalent, that is, true in exactly the same $\Sigma$-structures. This is the strongest reasonable reading of "the transformed claim preserves the original," and the paper's negative result concerns this strongest reading; weaker preservation relations (one-directional entailment, preservation of a designated consequence, and so on) inherit the same undecidability by essentially the same argument, since equivalence reduces to two entailment checks.

A \emph{transformation} is any map $T$ from claims to claims. No structure is assumed on $T$ in general; Section 4 restricts attention to transformations generated by a specific rewrite system.

\section{The General Correspondence Problem is Undecidable}

\begin{theorem}
There is no total algorithm $D : \mathrm{Sent}(\Sigma) \times \mathrm{Sent}(\Sigma) \to \{0,1\}$ such that $D(c,c') = 1$ if and only if $C(c,c')$, for all closed $\Sigma$-sentences $c, c'$.
\end{theorem}

\begin{proof}
Fix the sentence
\[
\tau := \forall x\,(P(x) \rightarrow P(x)),
\]
which is valid: it is satisfied by every $\Sigma$-structure regardless of the interpretation of $P$, $Q$, or $S$.

Let $\varphi$ range over closed $\Sigma$-sentences. Since $\tau$ is valid, every structure satisfying $\varphi$ also satisfies $\tau$, so $\varphi \models \tau$ holds unconditionally, for every $\varphi$. Consequently,
\[
C(\varphi, \tau) \iff (\varphi \models \tau) \wedge (\tau \models \varphi) \iff \tau \models \varphi.
\]
Since $\tau$ is valid, $\tau \models \varphi$ holds exactly when $\varphi$ is satisfied by every $\Sigma$-structure, that is, exactly when $\varphi$ is valid. Hence
\[
C(\varphi, \tau) \iff \varphi \text{ is valid}.
\]

Suppose, for contradiction, that a total decider $D$ for $C$ existed. Then $\varphi \mapsto D(\varphi, \tau)$ would decide validity for arbitrary closed $\Sigma$-sentences $\varphi$. But validity of first-order sentences over a signature containing a relation symbol of arity at least two is undecidable, by Church's theorem on the undecidability of first-order logic \cite{church1936}. This contradiction establishes that no total decider for $C$ exists.
\end{proof}

The requirement that $\Sigma$ contain a relation symbol of arity at least two is not a technicality. First-order logic restricted to signatures consisting only of unary predicates, the monadic fragment, is decidable \cite{lowenheim1915}. Church's undecidability result, and therefore Theorem 1, depends essentially on the presence of at least one relation of arity two or more; $S$ plays exactly this role in $\Sigma$, even though $S$ does not appear in $\tau$ itself. Only $\varphi$, ranging over the full language, needs to exploit $S$ for the reduction to inherit undecidability.

Theorem 1 says something narrower than "semantic correspondence between claims can never be established." It says that no single mechanical procedure decides correspondence correctly for every pair of claims expressible in a language this expressive. This rules out one particular and tempting design: a general-purpose $\mathrm{semanticEquivalent}()$ function, whether implemented by classical methods or by a learned model, that is assumed to give a correct verdict on an arbitrary pair of transformed and original claims. No such function can exist as a total, always-correct procedure. Any system that claims to certify correspondence in general is either incomplete, unsound, or restricted to a sub-language or a bounded class of transformations, whether or not it advertises that restriction.

\section{A Certified Fragment: Correspondence by Normalization}

Theorem 1 forecloses one design but does not foreclose all automation. If the class of admissible transformations is restricted rather than the claim language, correspondence can become decidable within that restricted class, and the restriction can still be expressive enough to matter.

\begin{definition}[Rewrite system]
A rewrite system $\mathcal{R}$ over $\mathrm{Sent}(\Sigma)$ is a set of rewrite rules $c \to c'$ together with the reduction relation $\to_{\mathcal{R}}$ they generate by application at any subformula, and its reflexive transitive closure $\to_{\mathcal{R}}^{*}$. $\mathcal{R}$ is \emph{terminating} if there is no infinite sequence $c_0 \to_{\mathcal{R}} c_1 \to_{\mathcal{R}} \cdots$, and \emph{locally confluent} if whenever $c \to_{\mathcal{R}} c_1$ and $c \to_{\mathcal{R}} c_2$, there exists $d$ with $c_1 \to_{\mathcal{R}}^{*} d$ and $c_2 \to_{\mathcal{R}}^{*} d$.
\end{definition}

By Newman's lemma \cite{newman1942}, a terminating and locally confluent rewrite system is confluent, and confluence together with termination guarantees that every sentence $c$ has a unique normal form $N_{\mathcal{R}}(c)$, the result of applying rules until none apply, independent of the order in which rules are applied.

\begin{definition}[Fragment correspondence]
Given a terminating, confluent rewrite system $\mathcal{R}$, define
\[
C_{\mathcal{R}}(c,c') \iff N_{\mathcal{R}}(c) = N_{\mathcal{R}}(c').
\]
\end{definition}

\begin{theorem}
For any terminating, confluent rewrite system $\mathcal{R}$ over $\mathrm{Sent}(\Sigma)$ with effectively computable rule application, $C_{\mathcal{R}}$ is decidable.
\end{theorem}

\begin{proof}
Termination guarantees that computing $N_{\mathcal{R}}(c)$ by repeated rule application halts after finitely many steps for any $c$. Confluence, obtained from termination and local confluence via Newman's lemma, guarantees that the resulting normal form is unique regardless of the order of rule application, so $N_{\mathcal{R}}(c)$ is a well-defined, effectively computable function of $c$. $C_{\mathcal{R}}(c,c')$ is then decided by computing $N_{\mathcal{R}}(c)$ and $N_{\mathcal{R}}(c')$ and testing syntactic equality, a decidable comparison on finite syntactic objects.
\end{proof}

The value of Theorem 2 depends entirely on $\mathcal{R}$ being non-trivial. A rewrite system whose only rules concern whitespace or byte-level formatting would make $C_{\mathcal{R}}$ decidable in a way that proves nothing of interest, since no realistic transformation would be mistaken for changing meaning under such a system in the first place. The following fragment is chosen to admit transformations that are lexically substantial while remaining semantically inert, so that deciding $C_{\mathcal{R}}$ by normalization does genuine work.

\begin{example}[A non-trivial certified fragment]
Let $\mathcal{R}$ consist of the following rule schemas, each preserving logical equivalence and each terminating and locally confluent under a fixed left-to-right orientation and a fixed canonical ordering on subformulas:
\begin{enumerate}
\item[(i)] alpha-renaming of bound variables to a canonical de Bruijn-indexed form;
\item[(ii)] reordering of conjuncts and disjuncts under a fixed total order on subformula syntax trees, exploiting associativity and commutativity of $\wedge$ and $\vee$;
\item[(iii)] double-negation elimination, $\neg\neg\psi \to \psi$;
\item[(iv)] elimination of redundant parenthesization and syntactic sugar with a single semantics-preserving unfolding, such as $\psi \rightarrow \chi \to (\neg \psi \vee \chi)$.
\end{enumerate}
Under this system,
\[
\forall x\,(P(x) \wedge Q(x)) \quad \text{and} \quad \forall y\,(Q(y) \wedge P(y))
\]
normalize to the same canonical form, via rule (i) renaming $y$ to the canonical bound variable and rule (ii) canonically ordering the conjuncts, despite being lexically distinct. This demonstrates that $\mathcal{R}$-certified correspondence is not merely byte identity performed after cosmetic preprocessing; it recognizes a genuine, if restricted, notion of semantic sameness.
\end{example}

The rewrite rules above are chosen with an eye toward a companion empirical paper: they form the natural benign side of a transformation battery applied to a real document, against which adversarial transformations, such as hypothesis deletion, quantifier alteration, or a changed constant, are deliberately placed outside $\mathcal{R}$'s domain.

\section{A Three-Valued Decision Procedure}

Theorem 1 and Theorem 2 together justify a decision procedure that is honest about the boundary between what it can certify and what it cannot.

\begin{definition}[Operational correspondence output]
Given a certified rewrite system $\mathcal{R}$ and a transformation $T$ taking $c$ to $c'$, define
Let $\mathrm{Dom}(\mathcal{R})$ denote the syntactic domain on which $\mathcal{R}$ is certified terminating and confluent, so $c,c'$ are said to \emph{lie in the certified fragment} when both belong to $\mathrm{Dom}(\mathcal{R})$ and $T$ is known to be composed of finitely many applications of rules in $\mathcal{R}$. Then
\[
\begin{small}
D_{\mathcal{R}}(c,c') =
\begin{cases}
\textsf{PRESERVED}, & \text{$c,c'$ in fragment, } N_{\mathcal{R}}(c) = N_{\mathcal{R}}(c'); \\[4pt]
\textsf{NOT\_PRESERVED\_UNDER\_}\mathcal{R}, & \text{$c,c'$ in fragment, } N_{\mathcal{R}}(c) \neq N_{\mathcal{R}}(c'); \\[4pt]
\textsf{OUTSIDE\_CERTIFIED\_FRAGMENT}, & \text{otherwise.}
\end{cases}
\end{small}
\]
\end{definition}

Two readings of this output deserve emphasis, because both are easy to collapse into a simpler but incorrect binary judgment.

First, $\textsf{NOT\_PRESERVED\_UNDER\_}\mathcal{R}$ does not entail $\neg C(c,c')$. It entails only that $\mathcal{R}$, the specific certified rewrite theory in use, did not establish correspondence. The two claims might still be logically equivalent for reasons outside $\mathcal{R}$'s rule set; the procedure simply has no certificate for that fact. Treating $\textsf{NOT\_PRESERVED\_UNDER\_}\mathcal{R}$ as proof of semantic change would silently smuggle back the false assumption that $\mathcal{R}$ is complete for correspondence, which Theorem 1 already rules out for any fixed mechanical procedure over the full language.

Second, $\textsf{OUTSIDE\_CERTIFIED\_FRAGMENT}$ is not a semantic verdict at all. It reports that the pair, or the transformation connecting them, falls outside the domain on which this mechanism claims automatic certification. It is deliberately distinct from an $\textsf{UNDECIDED}$ or $\textsf{UNKNOWN}$ output, which would suggest the procedure attempted to determine a truth value and failed. $\textsf{OUTSIDE\_CERTIFIED\_FRAGMENT}$ instead marks a boundary of applicability, and correctly routes the pair toward whatever fallback mechanism, formal reproof, symbolic checking, or attested human or expert judgment, a surrounding system chooses to require outside the certified regime.

The discipline this enforces can be stated compactly.

\begin{proposition}
Failure of certification is not certification of failure.
\end{proposition}

This proposition is not a further theorem so much as a naming of what Theorem 1 and the definition of $D_{\mathcal{R}}$ jointly imply: a bounded automatic procedure has exactly three honest outputs, a positive certificate, a negative certificate relative to its own rule set, and an admission of inapplicability, and none of the three may be read as a general semantic verdict beyond what the procedure actually established.

A companion protocol paper builds infrastructure around this boundary: a robust in-band commitment locating a signed claim manifest, transformation attestations recording which adjudication procedure $J$ was used and under what policy $P$ an adjudicator was authorized to use it, and inheritance of an earlier verification result conditioned on all of these. That architecture, informally the separation of locator recoverability, claim correspondence, adjudicator authorization, and verification inheritance, is best understood as the structure a system is forced into once $C$ cannot be silently collapsed into a universally computable step and adjudicator authority cannot be collapsed into "the verifier said so." Neither claim is defended here; both are consequences of taking Theorem 1 seriously and are developed in that companion paper. Correspondingly, whatever formal or informal adjudication regime $J$ is used outside the certified fragment, this paper's notation deliberately does not conflate a computed relation with an attested one: $c' \equiv_{\mathcal{L}} c$ denotes a formally defined and decided semantic relation within a language $\mathcal{L}$ that supports one, while $V : c' \approx_J c$ denotes a signed record of a judgment by verifier $V$ under regime $J$, which is evidence of a claim rather than a derivation of it.

\section{Related Work}

The undecidability argument in Theorem 1 rests on Church's classical result on the undecidability of first-order validity \cite{church1936}, and the necessity of a relation symbol of arity at least two follows from Löwenheim's decidability result for the monadic fragment \cite{lowenheim1915}. The decidability argument in Theorem 2 rests on Newman's lemma \cite{newman1942}, the standard route from termination and local confluence to unique normal forms in rewriting theory.

This paper's contribution is not a new watermarking or steganographic channel. Statistical text watermarking that alters token-sampling probabilities to leave a detectable trace is developed in SynthID-Text \cite{dathathri2024}, and the information-theoretic capacity of covert channels realized through unconstrained degrees of representational freedom is analyzed in the digital watermarking and steganography literature \cite{moulin2003, cachin2004}. This paper treats such mechanisms as infrastructure for transporting a robust locator, not as a contribution of its own, and its results are independent of which carrier technology is used.

Nor is this paper a new provenance ontology. Nanopublications formalize an atomic, signed scientific claim together with its provenance and publication metadata \cite{groth2010}, the W3C's PROV data model formalizes provenance relationships generally \cite{provo}, Verifiable Credentials formalize signed, verifiable assertions issued by one party about another \cite{vcdm}, and the Coalition for Content Provenance and Authenticity specifies machine-readable content history metadata \cite{c2pa}. Each of these supplies vocabulary and cryptographic machinery this paper's companion protocol paper can draw on directly; none of them addresses whether a claim survives a transformation with its meaning intact, which is the question this paper isolates.

The authorization question, which adjudicator is entitled to issue which class of judgment, is not new either. Trust management systems such as PolicyMaker \cite{blaze1996} and the SPKI/SDSI framework \cite{rivest1996, ellison1999} formalize exactly the separation between a principal being trusted to assert a predicate and the predicate's truth. This paper's companion protocol paper imports that lineage rather than proposing a new trust primitive; its contribution there is to prevent a correspondence judgment from acquiring authority merely because it is signed, using machinery this literature already supplies.

Two older traditions anticipate the structural problem this paper formalizes in a narrower setting. Proof repair, exemplified by systems that determine whether an existing formal proof survives a change to underlying definitions or libraries and, if not, what patch restores it \cite{ringer2020}, is the direct computational ancestor of the transformation-inheritance question posed here, restricted to the fully formal case. Citation-treatment systems in legal research, which track whether a cited authority is still followed, distinguished, criticized, or overruled rather than treating citation as a binary pointer, establish that "support inheritance after transformation" is a recurring problem across formal and institutional domains and is not specific to AI-generated or AI-transformed text.

\section{Conclusion}

This paper isolates a single question from the surrounding infrastructure of watermarking, provenance, and trust management: given a claim and a transformed successor, when may an earlier verification of the claim be automatically inherited by its successor? We show that the underlying correspondence relation, semantic equivalence between a claim and its transformed form, is undecidable in general once the claim language is expressive enough to encode first-order validity with a relation of arity at least two, and that it becomes decidable within a restricted, non-trivial class of transformations generated by a terminating and confluent rewrite system. The resulting decision procedure is deliberately three-valued rather than binary, distinguishing a positive certificate, a negative result relative to a specific rewrite theory, and an honest admission that a transformation falls outside what the procedure can certify at all.

The practical consequence is a design constraint rather than a system: any mechanism for verification inheritance under transformation, however it is built, must expose the boundary between the regime in which it can certify preservation and the regime in which it cannot, rather than collapsing every transformed pair into a single similarity judgment. A watermark can survive a transformation that changes a theorem's meaning; a signed attestation can be issued by an unauthorized party; a correspondence judgment can be attested without being derived. None of these failures is prevented by making the underlying carrier technology more robust. They are prevented, to the extent they can be prevented at all, by refusing to let persistence of a pointer stand in for validity of the claim it points to.

\begin{thebibliography}{99}

\bibitem{church1936}
A. Church, ``A note on the Entscheidungsproblem,'' \emph{Journal of Symbolic Logic}, 1936.

\bibitem{lowenheim1915}
L. L\"owenheim, ``\"Uber M\"oglichkeiten im Relativkalk\"ul,'' \emph{Mathematische Annalen}, 1915.

\bibitem{newman1942}
M. H. A. Newman, ``On theories with a combinatorial definition of `equivalence,''' \emph{Annals of Mathematics}, 1942.

\bibitem{dathathri2024}
S. Dathathri et al., ``Scalable watermarking for identifying large language model outputs,'' \emph{Nature}, 2024.

\bibitem{moulin2003}
P. Moulin and J. A. O'Sullivan, ``Information-theoretic analysis of information hiding,'' \emph{IEEE Transactions on Information Theory}, 2003.

\bibitem{cachin2004}
C. Cachin, ``An information-theoretic model for steganography,'' \emph{Information and Computation}, 2004.

\bibitem{groth2010}
P. Groth, A. Gibson, and J. Velterop, ``The anatomy of a nanopublication,'' \emph{Information Services and Use}, 2010.

\bibitem{provo}
World Wide Web Consortium, \emph{PROV-O: The PROV Ontology}, W3C Recommendation.

\bibitem{vcdm}
World Wide Web Consortium, \emph{Verifiable Credentials Data Model}, W3C Recommendation.

\bibitem{c2pa}
Coalition for Content Provenance and Authenticity, \emph{C2PA Technical Specification}.

\bibitem{blaze1996}
M. Blaze, J. Feigenbaum, and J. Lacy, ``Decentralized trust management,'' \emph{Proceedings of the IEEE Symposium on Security and Privacy}, 1996.

\bibitem{rivest1996}
R. L. Rivest and B. Lampson, ``SDSI: A Simple Distributed Security Infrastructure,'' 1996.

\bibitem{ellison1999}
C. Ellison et al., ``SPKI Certificate Theory,'' RFC 2693, 1999.

\bibitem{ringer2020}
T. Ringer et al., ``Proof repair across type equivalences,'' \emph{Proceedings of the ACM SIGPLAN Conference on Programming Language Design and Implementation (PLDI)}, 2020.

\end{thebibliography}

\end{document}
